\documentclass[12pt,a4paper]{report}
\usepackage[T1]{fontenc}
\usepackage[utf8]{inputenc}
\usepackage[polish,english]{babel}
\usepackage{lmodern}
\usepackage{amsmath,amssymb}
\usepackage{hyperref}
\usepackage{enumitem}
\usepackage{booktabs}
\usepackage{longtable}
\usepackage{graphicx}
\usepackage{listings}
\usepackage{xcolor}
\usepackage{tikz}
\usetikzlibrary{shapes,arrows,positioning}

\newcommand{\todo}[1]{\textcolor{red}{[TODO: #1]}}
\newcommand{\repo}{\texttt{AlphaPulse}}
\newcommand{\datasetver}{v5.2}
\begin{document}

\begin{titlepage}
\centering
{\Large Wydział Informatyki\par}
\vspace{0.3cm}
{\large Katedra Systemów Inteligentnych i Data Science\par}
\vspace{0.3cm}
{\large Inteligentne systemy przetwarzania danych\par}
\vspace{1.2cm}

{\large Kamil Czyżewski s22370\par}
{\large Kuba Szulc s23640\par}
\vspace{1cm}

{\LARGE Konfigurowalne środowisko AutoML dla finansowych potoków predykcyjnych w turnieju Numerai\par}
\vspace{1cm}

{\large Praca magisterska\par}
\vspace{1cm}

{\large Promotor techniczny Filip Dzięcioł\par}
{\large Promotor główny Grzegorz Wójcik\par}
\vfill
{\large Warszawa, Styczeń, 2026\par}
\end{titlepage}

\begin{titlepage}
\centering
{\Large Faculty of Computer Science\par}
\vspace{0.3cm}
{\large Department of Intelligent Systems and Data Science\par}
\vspace{1.2cm}

{\large Kamil Czyżewski s22370\par}
{\large Kuba Szulc s23640\par}
\vspace{1cm}

{\LARGE A Configurable AutoML Framework for Financial Prediction Pipelines in the Numerai Tournament\par}
\vspace{1cm}

{\large Master's thesis\par}
\vspace{1cm}

{\large Technical supervisor Filip Dzięcioł\par}
{\large Lead supervisor Grzegorz Wójcik\par}
\vfill
{\large Warsaw, January, 2026\par}
\end{titlepage}

\selectlanguage{polish}
\chapter*{Abstrakt}
Niniejsza praca magisterska przedstawia projekt i implementację systemu AlphaPulse --
konfigurowalnego frameworku do budowania, trenowania i wdrażania potoków uczenia
maszynowego przeznaczonych do udziału w turnieju Numerai. Numerai jest światowym
konkursem data science, w którym uczestniczący badacze dostarczają prognozy sygnałów
rynku akcji, które są agregowane następnie przez fundusz hedgingowy do realnych decyzji
inwestycyjnych.

Głównym problemem badawczym jest opracowanie abstrakcji i narzędzi, które umożliwiają
systematyczne eksperymentowanie z konfiguracjami modeli predykcyjnych dla danych
finansowych o charakterze szeregów czasowych -- przy zachowaniu rygorystycznych
wymagań dotyczących unikania wycieku informacji czasowej (look-ahead bias). W
odpowiedzi na ten problem zaproponowano framework oparty na deklaratywnej konfiguracji
YAML, hermetycznych abstrakcjach komponentów (preprocessory, modele, strategie
ensemble) oraz walidacji krzyżowej uwzględniającej strukturę er (purged era
cross-validation).

Główne wkłady pracy obejmują: modularną architekturę potoku umożliwiającą swobodne
komponowanie etapów przetwarzania cech i modeli, zautomatyzowany system optymalizacji
hiper parametrów (HPO) oparty na bibliotece Optuna, system AutoResearch -- autonomiczną
pętlę badawczą sterowaną agentowym modelem językowym, który iteracyjnie proponuje i
testuje modyfikacje konfiguracji potoku, oraz eksport predykcji zgodnych z wymaganiami
formatu submisji Numerai. Przeprowadzone eksperymenty potwierdzają, że system
umożliwia systematyczną eksplorację przestrzeni konfiguracji, a podejście agentowe
redukuje wymagany nakład pracy manualnej przy zachowaniu konkurencyjnej jakości
predykcji mierzonej metrykami Sharpe'a i średnią korelacją per-era.

\selectlanguage{english}
\chapter*{Abstract}
This Master's thesis describes the design and implementation of AlphaPulse -- a
configuration-driven system to build, train and deploy machine learning pipelines for the
Numerai tournament. Numerai is a worldwide data science challenge where users submit
stocks market predictions. Predictions that are compiled by a hedge fund and traded in
trades for real money.

The key research question is how to design abstractions and tooling to allow systematic
experimentation with predictive model configurations for financial time-series data, with a
strict temporal information leakage (look-ahead bias) prevention. Addressing this problem,
the thesis introduces its framework based on declarative YAML configuration, encapsulated
component abstractions (preprocessors, models, ensemble strategies) purged era cross
validation (and era-aware cross-validation).

A modular architecture for flexible composition of feature engineering stages using the
concept of pipeline. For each of them, they provide one or more models; an automated
hyperparameter optimization (HPO) system using the Optuna library; and the AutoResearch
system --- an autonomous research system. An agent-driven loop to propose and test
pipeline configuration mutations and an export pipeline for making predictions compatible
with the Numerai submission format.

The system shows systematic exploration of configuration spaces and the agent-driven
approach decreases the number of experiments performed, as demonstrated in conducted
experiments. In the end, maintaining the competitive prediction quality based on the per-era
Sharpe ratio and mean per-era correlation with a low manual research effort.

\selectlanguage{polish}
\chapter*{Słowa kluczowe}
uczenie maszynowe, prognozowanie rynku akcji, Numerai, optymalizacja
hiper parametrów, modele ensemble, duże modele językowe, agenci AI

\selectlanguage{english}
\chapter*{Keywords}
machine learning, stock market prediction, Numerai, hyperparameter
optimization, ensemble models, large language models, AI agents

\selectlanguage{polish}
\tableofcontents

\chapter*{Wstęp}

\section{Kontekst i motywacja}
Prognozowanie rynków finansowych należy do najtrudniejszych problemów uczenia
maszynowego. Dane finansowe wyróżniają się niskim stosunkiem sygnału do szumu,
niestacjonarnością, silną zależnością czasową między obserwacjami oraz strukturą grupową
(ery --- tygodniowe przekroje rynku), która wyklucza zastosowanie standardowych technik
podziału danych. Jednocześnie skuteczność predykcji jest weryfikowana przez rzeczywisty
rynek, co eliminuje ryzyko nadmiernego dopasowania do artefaktów zbioru testowego.

Turniej Numerai, uruchomiony w 2015 roku przez firmę Numerai Inc., stanowi unikalną
platformę badawczą w obszarze finansowego uczenia maszynowego. Numerai dostarcza
uczestnikom zaszyfrowane, zanonimizowane dane tabelaryczne opisujące spółki giełdowe i
ich atrybuty fundamentalne, techniczne oraz rynkowe. Zadaniem uczestników jest
wytrenowanie modeli predykujących przyszłe stopy zwrotu neutralne względem rynku
(alpha) w horyzoncie 20 dni roboczych. Prognozy wszystkich uczestników są agregowane
przez fundusz hedgingowy Numerai do jednego Meta Modelu, który steruje rzeczywistymi
decyzjami inwestycyjnymi. Uczestnicy mogą zbierać tokeny NMR (kryptowaluta Numerai) na
swoje modele, zdobywając nagrody za pozytywne wyniki lub tracąc środki w przypadku
słabych predykcji.

Ta architektura tworzy naturalny problem badawczy: jak systematycznie eksperymentować z
konfiguracjami modeli ML, aby maksymalizować jakość predykcji mierzoną metrykami takimi
jak korelacja per-era i współczynnik Sharpe'a, przy jednoczesnym spełnieniu
rygorystycznych wymagań dotyczących formatu submisji i unikania wycieku informacji
czasowej?

Istniejące podejścia do udziału w turnieju Numerai są w dużej mierze ad hoc: uczestnicy
piszą własne skrypty, eksperymentują manualnie i nie korzystają z ustandaryzowanych
abstrakcji. Brakuje gotowego, w pełni konfigurowalnego frameworku, który ujednoliciłby
cały przepływ pracy --- od pobierania danych przez definiowanie eksperymentów,
optymalizację hiper parametrów, walidację krzyżową z uwzględnieniem er, aż po eksport
pliku wynikowego w formacie wymaganym przez Numerai.

\section{Problem badawczy}
Niniejsza praca podejmuje następujący problem badawczy:

\begin{quote}
Jak zaprojektować i zaimplementować framework, który umożliwia systematyczne i
reprodukowalne eksperymentowanie z potokami uczenia maszynowego dla danych
finansowych turnieju Numerai, redukując jednocześnie nakład pracy inżynierskiej wymagany
do testowania nowych konfiguracji?
\end{quote}

Problem ten dekomponuje się na następujące pytania szczegółowe:
\begin{enumerate}[label=\arabic*.]
\item Jakie abstrakcje komponentów (preprocessory, modele, strategie ensemble) są
niezbędne do obsługi pełnego spektrum konfiguracji stosowanych w danych
finansowych o strukturze szeregów czasowych?
\item Jak zapewnić poprawną walidację krzyżową uwzględniającą zależności temporalne i
korelacje wewnątrz er, by uniknąć wycieku informacji?
\item W jakim stopniu automatyzacja poszukiwania hiperparametrów i konfiguracji modeli --
w tym z użyciem agentowego modelu językowego -- może zastąpić lub wspomóc
manualne eksperymenty badacza?
\item Jak zintegrować wszystkie etapy przepływu pracy w spójny, konfigurowalny system
zgodny z wymaganiami Numerai?
\end{enumerate}

\section{Cele pracy}
Głównym celem pracy jest zaprojektowanie i zaimplementowanie frameworku AlphaPulse,
który:
\begin{itemize}
\item Ujednolica pełny przepływ pracy konkursu Numerai --- od pobierania danych po eksport submisji.
\item Umożliwia deklaratywne definiowanie eksperymentów w formacie YAML, bez
konieczności modyfikowania kodu źródłowego.
\item Zapewnia poprawną walidację z uwzględnieniem struktury er, purging i embargo.
\item Automatyzuje poszukiwanie optymalnych hiperparametrów i konfiguracji pipeline'u za
pomocą Optuna.
\item Wprowadza autonomiczną pętlę badawczą (AutoResearch) sterowaną dużym
modelem językowym (LLM), która iteracyjnie proponuje i testuje modyfikacje
konfiguracji.
\item Eksportuje wyniki w postaci pliku \texttt{predict.pkl} zgodnego z wymaganiami API Numerai.
\end{itemize}

Celami pomocniczymi są: ocena skuteczności podejścia agentowego w kontekście AutoML,
analiza wpływu wyboru modeli i strategii ensemble na metryki jakości predykcji oraz
porównanie różnych strategii walidacji krzyżowej dla danych finansowych z strukturą er.

\section{Zakres i ograniczenia}
Framework AlphaPulse jest przeznaczony wyłącznie dla turnieju Numerai Classic i bazuje
na zestawie danych w wersji 5.2. Zakres pracy obejmuje:
\begin{itemize}
\item Modele tabelaryczne: gradient boosting, modele drzewiste, modele liniowe oraz
fundamentalne modele tabelaryczne.
\item Preprocessory: skalowanie, redukcja wymiarowości, selekcja cech, regularyzacja
poprzez dodanie szumu gaussowskiego do danych.
\item Strategie ensemble: weighting ensemble, stacking z meta-modelem.
\item Optymalizację hiperparametrów z użyciem biblioteki Optuna.
\item System AutoResearch na podstawie frameworku Andreja Karpathy’ego z agentem
badawczym.
\end{itemize}

Praca nie obejmuje: analizy modeli sekwencyjnych, czyli na przykład LSTM czy Transformer,
dla danych czasowych, strategii tradingowych ani mechanizmów stakowania tokenów NMR,
a także integracji z turniejem Numerai Signals.

\section{Wkład autorów}
Wkład pracy obejmuje:
\begin{itemize}
    \item projekt i implementację frameworku \repo{} (Python 3.12+, branch \texttt{feat/core-framework}),
    \item deklaratywny format eksperymentów (Experiment v1 YAML),
    \item integrację HPO (Optuna) oraz pętli agentowej AutoResearch,
    \item ujednolicony pipeline od pobrania danych Numerai do eksportu \texttt{predict.pkl}.
\end{itemize}
\section{Kryteria sukcesu}
Praca uznaje się za spełnioną, jeśli:
\begin{enumerate}
    \item framework umożliwia powtarzalne uruchomienia tej samej konfiguracji,
    \item HPO poprawia metrykę główną względem baseline o co najmniej \todo{wpisz $\Delta$ corr},
    \item AutoResearch redukuje liczbę manualnych iteracji o co najmniej \todo{wpisz \%},
    \item końcowy artefakt submisji jest zgodny z API Numerai.
\end{enumerate}

\chapter{Podstawy problemu i kontekst systemu AlphaPulse}

\section{Cel rozdziału}
Celem niniejszego rozdziału jest przedstawienie kontekstu domenowego i technicznego,
w którym powstał framework \textit{AlphaPulse}. Rozdział porządkuje pojęcia wymagane do
zrozumienia dalszych części pracy: specyfikę turnieju Numerai, charakter danych,
kryteria oceny modeli, ograniczenia metodologiczne oraz wymagania, jakie musi spełniać
nowoczesny system AutoML dla finansowych potoków predykcyjnych.

\section{Specyfika turnieju Numerai}
Numerai jest konkursem data science, w którym uczestnicy budują modele predykcyjne
na zanonimizowanych danych tabelarycznych i dostarczają predykcje do wspólnego
meta-modelu funduszu. W przeciwieństwie do klasycznych konkursów ML, nacisk jest
położony nie tylko na jakość predykcji, ale również na stabilność sygnału i odporność
na zmianę warunków rynkowych.

W praktyce oznacza to, że:
\begin{itemize}
    \item dane mają strukturę czasową i grupową (ery),
    \item błędna walidacja może prowadzić do sztucznie zawyżonych wyników,
    \item wysokie znaczenie ma reprodukowalność eksperymentów,
    \item końcowym artefaktem jest plik predykcji zgodny z wymaganym formatem submisji.
\end{itemize}

Takie środowisko premiuje podejście systemowe: powtarzalne pipeline'y, kontrolę konfiguracji,
jasne logowanie eksperymentów i automatyzację przeszukiwania przestrzeni modeli.

\section{Charakterystyka danych finansowych}
Dane używane w zadaniach Numerai należą do kategorii danych tabelarycznych o silnym
komponencie czasowym. Z perspektywy modelowania najważniejsze cechy tych danych to:
\begin{itemize}
    \item niski stosunek sygnału do szumu,
    \item niestacjonarność rozkładów cech i celu,
    \item zależności temporalne między obserwacjami,
    \item podatność na przeciek informacji przy niepoprawnym podziale danych.
\end{itemize}

Konsekwencją jest konieczność ostrożnego projektowania zarówno preprocessingu,
jak i procedury walidacyjnej. W praktyce szczególnie istotne stają się metody walidacji
uwzględniające strukturę er oraz separację czasową między zbiorem uczącym i walidacyjnym.

\section{Metryki jakości i kryteria oceny}
W odróżnieniu od wielu zadań regresyjnych, ocena modeli finansowych nie powinna
opierać się wyłącznie na błędzie punktowym. W kontekście Numerai kluczowe są metryki
oceniające jakość sygnału per-era oraz stabilność tego sygnału w czasie.

W pracy przyjmuje się dwa podstawowe kryteria:
\begin{itemize}
    \item średnia korelacja per-era, jako miara jakości predykcji względem celu,
    \item współczynnik Sharpe'a liczony na poziomie sygnału per-era, jako miara relacji
    jakości do zmienności.
\end{itemize}

Takie ujęcie premiuje modele nie tylko skuteczne, lecz także stabilne, co ma bezpośredni
związek z użytecznością predykcji w zastosowaniu konkursowym.

\section{Ryzyko wycieku informacji i poprawna walidacja}
Jednym z najważniejszych problemów metodologicznych w modelach finansowych jest
\textit{look-ahead bias}, czyli niejawne wykorzystanie informacji, która w momencie
predykcji nie byłaby dostępna. Zjawisko to może wystąpić na wielu etapach:
\begin{itemize}
    \item podczas przygotowania cech,
    \item przy podziale danych na train/validation,
    \item przy doborze hiperparametrów.
\end{itemize}

Dlatego system klasy \textit{AlphaPulse} musi wspierać walidację zgodną z naturą danych
czasowych i erowych, w tym podejścia \textit{era-aware} oraz \textit{purged cross-validation}.
Wymóg ten stanowi fundament wiarygodnej oceny i bezpośrednio wpływa na konstrukcję
architektury frameworku.

\section{Od eksperymentów ad hoc do frameworku konfigurowalnego}
W praktyce uczestnicy konkursów często rozwijają rozwiązania jako zbiór luźnych skryptów.
Taki model pracy utrudnia:
\begin{itemize}
    \item odtwarzanie wyników,
    \item porównywanie wariantów modeli,
    \item skalowanie liczby eksperymentów,
    \item bezpieczne wdrażanie zmian.
\end{itemize}

Odpowiedzią na ten problem jest podejście konfiguracyjne (\textit{configuration-driven}),
w którym eksperyment opisuje się deklaratywnie, a kod frameworku odpowiada za
deterministyczne wykonanie potoku. Właśnie ten paradygmat został przyjęty w projekcie
\textit{AlphaPulse}.

\section{Zakres funkcjonalny AlphaPulse jako punkt odniesienia dalszych rozdziałów}
Na poziomie operacyjnym framework obejmuje pełny przepływ pracy konkursowej:
\begin{enumerate}
    \item pobranie danych,
    \item uruchamianie eksperymentów opisanych YAML,
    \item automatyczną optymalizację hiperparametrów (HPO),
    \item agentową pętlę badawczą (\textit{AutoResearch}),
    \item eksport predykcji do formatu submisji.
\end{enumerate}

W projekcie uwzględniono również:
\begin{itemize}
    \item modularny preprocessing (m.in. skalowanie, redukcja wymiarowości, selekcja cech),
    \item szeroką rodzinę modeli tabelarycznych (boosting, modele drzewiaste, liniowe,
    wybrane modele fundamentowe dla danych tabelarycznych),
    \item strategie ensemble (single, weighted, stacking),
    \item standaryzację uruchomień przez zestaw skryptów eksperymentalnych.
\end{itemize}

Ten zakres stanowi podstawę układu kolejnych rozdziałów: najpierw architektura,
następnie implementacja i konfiguracja eksperymentów, a dalej ewaluacja wyników.

\section{Mapowanie wyzwań domenowych na decyzje projektowe}
\centering
\begin{longtable}{p{4.2cm}p{4.2cm}p{5.2cm}}
\caption{Wyzwanie finansowe $\rightarrow$ decyzja w \repo{}}\\
\toprule
Wyzwanie & Ryzyko & Decyzja projektowa \\
\midrule
Struktura er & błędny split & walidacja era-aware, purge/embargo \\
Niestacjonarność & niestabilny sygnał & metryki per-era + Sharpe \\
Duża przestrzeń konfiguracji & koszt manualny & HPO + AutoResearch \\
Brak standaryzacji workflow & brak reprodukowalności & YAML + CLI + artefakty \\
Wymóg submisji & błąd formatu & dedykowany eksport \texttt{predict.pkl} \\
\bottomrule
\end{longtable}


\section{Wnioski}
Przedstawiony kontekst pokazuje, że skuteczne modelowanie w Numerai wymaga połączenia
trzech perspektyw: poprawności metodologicznej, automatyzacji eksperymentów oraz
inżynierii oprogramowania zapewniającej reprodukowalność. Z tego powodu w dalszej
części pracy \textit{AlphaPulse} jest traktowany nie jako pojedynczy model, lecz jako
system badawczo-produkcyjny wspierający iteracyjne budowanie i ocenę pipeline'ów ML.

\chapter{Architektura systemu AlphaPulse}

\section{Cel rozdziału}
Celem rozdziału jest przedstawienie architektury frameworku \textit{AlphaPulse} jako
konfigurowalnego systemu AutoML dla danych finansowych Numerai. Opis obejmuje
podział na warstwy, przepływ danych, główne komponenty wykonawcze oraz sposób,
w jaki architektura wspiera reprodukowalność i iteracyjny rozwój eksperymentów.

\section{Założenia architektoniczne}
Projekt \textit{AlphaPulse} został oparty na kilku kluczowych założeniach:
\begin{itemize}
    \item \textbf{Konfigurowalność}: eksperymenty są definiowane deklaratywnie, bez
    każdorazowej modyfikacji kodu źródłowego.
    \item \textbf{Modularność}: preprocessing, modele i strategie ensemble są odseparowane
    i mogą być dowolnie komponowane.
    \item \textbf{Reprodukowalność}: uruchomienia mają deterministyczny charakter
    przy ustalonej konfiguracji i ziarnie losowym.
    \item \textbf{Rozszerzalność}: nowe komponenty mogą być dodawane bez naruszania
    istniejącego przepływu.
    \item \textbf{Automatyzacja}: architektura wspiera zarówno klasyczne eksperymenty,
    jak i zautomatyzowane przeszukiwanie konfiguracji (HPO, AutoResearch).
\end{itemize}

\section{Widok warstwowy systemu}
Architekturę systemu można opisać jako układ pięciu współpracujących warstw.

\subsection{Warstwa danych}
Warstwa danych odpowiada za:
\begin{itemize}
    \item pobranie i lokalne utrzymanie zbiorów Numerai,
    \item ładowanie podziałów (np. train/validation),
    \item przygotowanie macierzy cech, celu i metadanych (era),
    \item kontrolę parametrów wejściowych (np. subsampling, feature set).
\end{itemize}

Ta warstwa dostarcza ustandaryzowany interfejs wejściowy dla całego pipeline'u.

\subsection{Warstwa konfiguracji eksperymentu}
Konfiguracja eksperymentu jest deklarowana w formacie YAML (schema-driven).
Opisuje m.in.:
\begin{itemize}
    \item źródło danych i parametry podziału,
    \item zestawy cech i grupy cech,
    \item listę preprocessingu,
    \item typy modeli oraz ich hiperparametry,
    \item strategię ensemble,
    \item parametry treningu i ewaluacji.
\end{itemize}

Dzięki temu logika eksperymentu jest jawna i łatwa do wersjonowania.

\subsection{Warstwa pipeline'u modelowego}
Warstwa pipeline'u realizuje sekwencję:
\[
\text{Data} \rightarrow \text{Preprocessing} \rightarrow \text{Model(e)} \rightarrow \text{Ensemble} \rightarrow \text{Prediction}
\]
W ramach tej warstwy:
\begin{itemize}
    \item preprocessory przekształcają cechy wejściowe,
    \item modele bazowe uczą się predykcji celu,
    \item ensemble agreguje sygnały modeli cząstkowych.
\end{itemize}

Pipeline jest budowany dynamicznie na podstawie konfiguracji, co umożliwia szybkie
porównywanie wielu wariantów architektury modelowej.

\subsection{Warstwa optymalizacji i automatyzacji}
System wspiera dwa poziomy automatyzacji:
\begin{itemize}
    \item \textbf{HPO (Optuna)}: automatyczne przeszukiwanie przestrzeni hiperparametrów
    i konfiguracji komponentów.
    \item \textbf{AutoResearch}: agentowa pętla badawcza, która iteracyjnie proponuje
    mutacje konfiguracji i ocenia ich wpływ na metryki.
\end{itemize}

Wynikiem działania tej warstwy jest ranking prób, najlepsza konfiguracja oraz historia
decyzji eksperymentalnych.

\subsection{Warstwa artefaktów i eksportu}
Warstwa artefaktów odpowiada za:
\begin{itemize}
    \item zapisy metryk i logów eksperymentów,
    \item serializację najlepszych konfiguracji,
    \item eksport predykcji do formatu wymaganej submisji (np. \texttt{predict.pkl}).
\end{itemize}

To domyka przepływ od eksperymentu badawczego do artefaktu konkursowego.

\section{Przepływ wykonania end-to-end}
Typowy scenariusz działania systemu obejmuje następujące etapy:
\begin{enumerate}
    \item pobranie danych i przygotowanie lokalnego katalogu roboczego,
    \item uruchomienie eksperymentu z pliku YAML,
    \item trenowanie i ewaluacja modeli bazowych oraz ensemble,
    \item (opcjonalnie) uruchomienie HPO lub AutoResearch,
    \item wybór najlepszej konfiguracji i eksport predykcji.
\end{enumerate}

W ujęciu operacyjnym odpowiada temu zestaw dedykowanych skryptów CLI, które
oddzielają definicję eksperymentu od mechaniki wykonania.

\section{Komponenty funkcjonalne i ich odpowiedzialności}
Architektura dzieli odpowiedzialności pomiędzy komponenty:
\begin{itemize}
    \item \textbf{Data Loader}: odczyt danych i przygotowanie splitów,
    \item \textbf{Experiment Schema}: walidacja konfiguracji wejściowej,
    \item \textbf{Pipeline Builder}: konstrukcja grafu preprocessingu, modeli i ensemble,
    \item \textbf{Trainer/Evaluator}: trening i obliczanie metryk jakości,
    \item \textbf{HPO Engine}: zarządzanie próbami optymalizacyjnymi,
    \item \textbf{AutoResearch Loop}: agentowe sterowanie kolejnymi eksperymentami,
    \item \textbf{Exporter}: generowanie artefaktów konkursowych.
\end{itemize}

Takie rozdzielenie upraszcza testowanie i ułatwia rozwój systemu.

\section{Wzorce zapewniające jakość inżynierską}
W \textit{AlphaPulse} zastosowano praktyki zwiększające niezawodność:
\begin{itemize}
    \item \textbf{Separation of concerns} między konfiguracją, logiką uczenia i eksportem,
    \item \textbf{CLI-first workflow} dla łatwej automatyzacji uruchomień,
    \item \textbf{Jednoznaczne artefakty wyjściowe} dla porównywania eksperymentów,
    \item \textbf{Integracja testów i narzędzi jakości kodu} dla stabilności rozwoju.
\end{itemize}

W efekcie architektura nadaje się zarówno do badań iteracyjnych, jak i do utrzymania
dłuższej linii rozwojowej projektu.

\section{Ograniczenia architektury}
Pomimo wysokiej elastyczności architektura ma ograniczenia:
\begin{itemize}
    \item skuteczność zależy od jakości zdefiniowanej przestrzeni konfiguracji,
    \item duża liczba wariantów zwiększa koszt obliczeniowy eksperymentów,
    \item automatyzacja agentowa wymaga kontroli jakości proponowanych mutacji.
\end{itemize}

Ograniczenia te są naturalne dla systemów AutoML i stanowią punkt wyjścia do
dalszej optymalizacji.

\section{Diagram architektury logicznej}
\begin{figure}[h]
\centering
\begin{tikzpicture}[node distance=1.4cm, auto]
    \node (data) [draw, rectangle] {Warstwa danych};
    \node (cfg) [draw, rectangle, below of=data] {Konfiguracja YAML};
    \node (pipe) [draw, rectangle, below of=cfg] {Pipeline ML};
    \node (opt) [draw, rectangle, below of=pipe] {HPO / AutoResearch};
    \node (out) [draw, rectangle, below of=opt] {Artefakty / Export};
    \draw[->] (data) -- (cfg);
    \draw[->] (cfg) -- (pipe);
    \draw[->] (pipe) -- (opt);
    \draw[->] (opt) -- (out);
\end{tikzpicture}
\caption{Przepływ end-to-end w \repo{}}
\end{figure}

\section{Kontrakty komponentów}
Kontrakty komponentów określają minimalne wymagania wejścia i wyjścia dla kolejnych
warstw systemu. Dzięki temu pipeline może być rozwijany modularnie: nowy loader,
preprocessor, model lub strategia ensemble musi spełnić ten sam interfejs logiczny, bez
konieczności zmiany pozostałych części frameworku.

\begin{longtable}{p{3.2cm}p{5.2cm}p{5.2cm}}
\caption{Podstawowe kontrakty komponentów systemu \repo{}}\\
\toprule
Komponent & Wejście & Wyjście \\
\midrule
Loader danych & konfiguracja zbioru Numerai, wersja danych, lista cech & macierze \texttt{X}, \texttt{y}, kolumna \texttt{era} oraz metadane zbioru \\
Preprocessor & dane treningowe i parametry transformacji & przekształcone cechy oraz zapisany stan transformacji \\
Model bazowy & cechy, target, konfiguracja hiperparametrów & wytrenowany estymator i predykcje walidacyjne \\
Ensemble & predykcje modeli bazowych i konfiguracja łączenia & finalna predykcja oraz metadane wag lub meta-modelu \\
Eksport & predykcje końcowe i identyfikatory przykładów & plik \texttt{predict.pkl} zgodny z wymaganiami Numerai \\
\bottomrule
\end{longtable}

Konfiguracja eksperymentu pełni rolę wspólnego kontraktu sterującego. Wymagane sekcje
YAML obejmują \texttt{data}, \texttt{features}, \texttt{preprocessing}, \texttt{models},
\texttt{ensemble\_method}, \texttt{train} oraz \texttt{evaluation}. Każde uruchomienie
zapisuje konfigurację, metryki, logi prób oraz, zależnie od trybu pracy, artefakty takie jak
\texttt{best\_config.json} i \texttt{predict.pkl}.

\section{Podsumowanie}
Architektura \textit{AlphaPulse} realizuje podejście \textit{configuration-driven} i integruje
pełny cykl pracy konkursowej: od danych, przez eksperymenty i automatyzację, po eksport
predykcji. Kluczową wartością systemu jest połączenie modularności technicznej z
reprodukowalnością badawczą, co stanowi fundament dalszych rozdziałów implementacyjnych
i eksperymentalnych.

\chapter{Implementacja systemu AlphaPulse}

\section{Cel rozdziału}
Celem rozdziału jest szczegółowe przedstawienie implementacji frameworku
\textit{AlphaPulse}: sposobu organizacji kodu, mechanizmów konfiguracji eksperymentów,
interfejsów uruchomieniowych oraz realizacji kluczowych etapów pipeline'u od danych do
artefaktu submisyjnego.

\section{Organizacja repozytorium}
Implementacja projektu została podzielona na warstwy odpowiadające cyklowi pracy
eksperymentalnej:
\begin{itemize}
    \item \texttt{src/alphapulse/} -- kod domenowy i komponenty frameworku,
    \item \texttt{scripts/} -- skrypty wykonawcze CLI dla najważniejszych scenariuszy,
    \item \texttt{experiments/} -- deklaratywne definicje eksperymentów (YAML),
    \item \texttt{tests/} -- testy jednostkowe i integracyjne,
    \item \texttt{artifacts/} -- wyniki uruchomień, konfiguracje i eksporty.
\end{itemize}

Taki podział rozdziela logikę frameworku od warstwy operacyjnej (uruchamianie),
co zwiększa czytelność i ułatwia utrzymanie projektu.

\section{Warstwa uruchomieniowa (CLI)}
\subsection{Skrypty wykonawcze}
System udostępnia zestaw skryptów CLI obsługujących pełny przepływ:
\begin{itemize}
    \item pobieranie danych,
    \item uruchamianie eksperymentów YAML,
    \item HPO,
    \item lekki test end-to-end (\textit{test pipeline}),
    \item eksport predykcji do formatu Numerai.
\end{itemize}

Skrypty pełnią rolę cienkiej warstwy orkiestracji: pobierają argumenty, walidują wejście,
inicjują komponenty frameworku i zapisują artefakty.

\subsection{Konwencja parametrów wejściowych}
Interfejs CLI wykorzystuje jednolitą konwencję argumentów:
\begin{itemize}
    \item ścieżki danych i artefaktów (\texttt{--data-dir}, \texttt{--output-dir}),
    \item parametry skali eksperymentu (\texttt{--train-subsample}, liczba prób),
    \item wskazanie konfiguracji źródłowej (\texttt{--config}, \texttt{--best-config-path}),
    \item tryb lokalny lub rozproszony (zależnie od skryptu).
\end{itemize}

Ujednolicenie argumentów upraszcza automatyzację i integrację z narzędziami CI/CD.

\section{Konfiguracja eksperymentów (YAML)}
\subsection{Schemat i walidacja}
Eksperymenty definiowane są deklaratywnie. Schemat obejmuje sekcje:
\begin{itemize}
    \item \texttt{data},
    \item \texttt{features},
    \item \texttt{preprocessing},
    \item \texttt{models},
    \item \texttt{ensemble\_method}/\texttt{ensemble\_params},
    \item \texttt{train},
    \item \texttt{evaluation}.
\end{itemize}

Każde uruchomienie przechodzi walidację zgodności konfiguracji, co ogranicza liczbę
błędów wykonania i gwarantuje spójność eksperymentów.

\subsection{Konfigurowalność bez zmian kodu}
Najważniejszą własnością implementacji jest możliwość modyfikacji eksperymentu
poprzez edycję YAML:
\begin{itemize}
    \item zmiana modelu i hiperparametrów,
    \item podmiana preprocessingu,
    \item przełączenie strategii ensemble,
    \item zmiana metryki głównej ewaluacji.
\end{itemize}

Pozwala to prowadzić szybkie iteracje badawcze przy zachowaniu pełnej ścieżki audytu
zmian w repozytorium.

\section{Implementacja warstwy danych}
\subsection{Loader danych Numerai}
Warstwa danych dostarcza zunifikowany interfejs ładowania splitów oraz metadanych.
Implementacja obejmuje:
\begin{itemize}
    \item odczyt plików parquet i plików opisowych,
    \item wybór zestawu cech i subsamplingu,
    \item przygotowanie struktur \texttt{X}, \texttt{y}, \texttt{era}.
\end{itemize}

Taki kontrakt danych jest wykorzystywany przez wszystkie wyższe warstwy pipeline'u.

\subsection{Obsługa danych wejściowych pod eksperymenty}
Loader obsługuje zarówno pełne zbiory, jak i tryby szybkich testów na podzbiorach danych.
To umożliwia:
\begin{itemize}
    \item szybkie iteracje deweloperskie,
    \item tańsze testy poprawności potoku,
    \item stopniowe skalowanie kosztownych eksperymentów.
\end{itemize}

\section{Implementacja pipeline'u modelowego}
\subsection{Preprocessing}
Implementacja preprocessingu wykorzystuje listę kroków konfigurowanych w YAML.
Wspierane są m.in.:
\begin{itemize}
    \item skalowanie (\textit{StandardScaler}, \textit{RobustScaler}),
    \item redukcja wymiarowości (\textit{PCA}, \textit{TruncatedSVD}),
    \item selekcja cech i transformacje regularizujące.
\end{itemize}

Kroki preprocessingu są wykonywane deterministycznie i mogą być łączone w sekwencje.

\subsection{Modele bazowe}
Warstwa modeli została zaprojektowana jako zestaw wymiennych adapterów
(spójny interfejs \textit{fit/predict}) dla różnych rodzin:
\begin{itemize}
    \item gradient boosting,
    \item modele drzewiaste,
    \item modele liniowe,
    \item wybrane modele fundamentowe dla danych tabelarycznych.
\end{itemize}

Standaryzacja interfejsu upraszcza porównywanie modeli i ich integrację z ensemble.

\subsection{Strategie ensemble}
Implementacja obsługuje trzy podstawowe tryby:
\begin{itemize}
    \item \texttt{single} -- pojedynczy model,
    \item \texttt{weighted} -- ważona agregacja predykcji,
    \item \texttt{stacking} -- meta-model uczony na predykcjach bazowych.
\end{itemize}

Strategia ensemble jest traktowana jako pełnoprawny element konfiguracji eksperymentu.

\section{Implementacja HPO}
\subsection{Integracja z Optuna}
HPO realizowane jest przez warstwę opartą o Optuna, która:
\begin{itemize}
    \item definiuje przestrzeń przeszukiwania dla hiperparametrów i wariantów pipeline'u,
    \item uruchamia próby treningowo-ewaluacyjne,
    \item rankinguje konfiguracje względem metryki celu.
\end{itemize}

Wynikiem jest najlepsza konfiguracja zapisywana jako artefakt możliwy do dalszego
użycia (np. eksport predykcji).

\subsection{Śledzenie wyników prób}
Każda próba HPO generuje komplet informacji eksperymentalnych:
\begin{itemize}
    \item parametry wejściowe,
    \item metryki jakości,
    \item metadane wykonania.
\end{itemize}
Umożliwia to analizę post hoc i budowę wiedzy o wrażliwości systemu na konfigurację.

\section{Implementacja AutoResearch}
\subsection{Pętla agentowa}
\textit{AutoResearch} rozszerza klasyczny HPO o mechanizm agentowy:
\begin{enumerate}
    \item analiza dotychczasowych wyników,
    \item propozycja mutacji konfiguracji,
    \item uruchomienie kolejnej próby,
    \item aktualizacja stanu badania.
\end{enumerate}

Stan eksperymentu i decyzje agenta są zapisywane w artefaktach, co pozwala audytować
przebieg procesu badawczego.

\subsection{Rola AutoResearch w implementacji}
Mechanizm ten nie zastępuje pełni pracy badacza, lecz automatyzuje najbardziej
powtarzalny etap: generowanie kolejnych hipotez konfiguracyjnych i ich testowanie.
W praktyce skraca to czas dochodzenia do konkurencyjnych wariantów pipeline'u.

\section{Eksport predykcji i artefakty końcowe}
Końcowym etapem implementacji jest produkcja artefaktu submisyjnego:
\begin{itemize}
    \item odtworzenie najlepszego pipeline'u,
    \item trening na wskazanym zbiorze,
    \item wygenerowanie predykcji,
    \item serializacja do formatu kompatybilnego z Numerai (np. \texttt{predict.pkl}).
\end{itemize}

Spójny format wyjścia domyka proces od konfiguracji badawczej do użycia konkursowego.

\section{Zapewnienie jakości kodu}
W implementacji zastosowano standardowe mechanizmy jakości:
\begin{itemize}
    \item linting i formatowanie kodu,
    \item statyczne sprawdzanie typów,
    \item testy automatyczne,
    \item hooki pre-commit i praktyki CI.
\end{itemize}

To ogranicza regresje i zwiększa stabilność frameworku w kolejnych iteracjach.

\section{Katalog komponentów wspieranych przez system}
\label{sec:katalog-komponentow}

\subsection{Modele bazowe}
Framework wspiera następujące typy modeli:
\begin{itemize}
    \item \textbf{Gradient boosting:} XGBoost, LightGBM, CatBoost, Packboost,
    \item \textbf{Ensemble drzewiaste:} RandomForest, ExtraTrees,
    \item \textbf{Modele liniowe:} Ridge,
    \item \textbf{Foundation models:} TabPFN, TabICL,
    \item \textbf{Meta-modele:} EraEnsemble, SyntheticDataAugmenter.
\end{itemize}

\subsection{Preprocessory}
\begin{itemize}
    \item skalowanie: StandardScaler, RobustScaler,
    \item redukcja wymiarowości: PCA, TruncatedSVD,
    \item selekcja/regularizacja: VarianceSelector, LGBMImportanceSelector, GaussianNoise,
    \item specjalistyczne: Packboost, GroupedPreprocessor.
\end{itemize}

\subsection{Strategie ensemble}
\begin{itemize}
    \item \texttt{single} -- jeden model,
    \item \texttt{weighted} -- agregacja ważona,
    \item \texttt{stacking} -- meta-learner (\texttt{ridge} lub \texttt{xgboost}).
\end{itemize}

\section{Przykładowa konfiguracja referencyjna}
\begin{lstlisting}[language=Python,basicstyle=\small\ttfamily]
version: "1"
data:
  data_dir: data/v5.2
  train_subsample: 0.05
  target_col: target
  seed: 42
models:
  - type: XGBoost
    params:
      max_depth: 3
      learning_rate: 0.05
      tree_method: hist
ensemble_method: single
evaluation:
  primary_metric: mean_per_era_correlation
\end{lstlisting}

\section{Przestrzeń przeszukiwania HPO}
Przykładowe parametry optymalizowane przez Optuna:
\begin{itemize}
    \item typ modelu (XGBoost/LightGBM/CatBoost/Ridge/...),
    \item hiperparametry modelu (\texttt{max\_depth}, \texttt{learning\_rate}, \texttt{n\_estimators}),
    \item wybór preprocessingu i jego parametrów,
    \item metoda ensemble i parametry meta-modelu.
\end{itemize}
Budżet domyślny: \todo{np. 30 prób HPO, 50 prób AutoResearch, 2h czasu}.

\section{Mutacje AutoResearch}
Agent może proponować:
\begin{enumerate}
    \item dodanie/usunięcie modelu bazowego,
    \item zmianę preprocessingu,
    \item zmianę strategii ensemble,
    \item strojenie hiperparametrów,
    \item ustawienie neutralizacji cech (jeśli aktywne w konfiguracji).
\end{enumerate}
Stan badania zapisywany jest w \texttt{research\_state.json}, a podsumowanie prób w \texttt{trials\_summary.csv}.

\section{Podsumowanie}
Implementacja \textit{AlphaPulse} realizuje wymagania postawione we wcześniejszych
rozdziałach: modularność, konfigurowalność i reprodukowalność. Dzięki warstwie CLI,
schematom YAML, integracji HPO oraz pętli AutoResearch system wspiera zarówno szybkie
eksperymenty, jak i systematyczne badania nad potokami predykcyjnymi dla Numerai.

\chapter{Metodyka badań i plan eksperymentów}

\section{Cel rozdziału}
Celem rozdziału jest zdefiniowanie metodyki oceny frameworku \textit{AlphaPulse} oraz
przedstawienie planu eksperymentów umożliwiającego rzetelną odpowiedź na pytania
badawcze postawione we wstępie. Rozdział formalizuje: cele pomiaru, układ porównań,
protokół uruchomień, metryki, sposób analizy wyników oraz kryteria interpretacji.

\section{Pytania badawcze i hipotezy}
Na podstawie problemu badawczego przyjęto następujące pytania operacyjne:

\begin{enumerate}
    \item Czy konfiguracja deklaratywna i modularna architektura pipeline'u zwiększają
    reprodukowalność oraz porównywalność eksperymentów?
    \item Jaki wpływ na jakość predykcji mają: preprocessing, wybór modelu bazowego
    oraz strategia ensemble?
    \item W jakim stopniu HPO poprawia metryki jakości względem konfiguracji bazowych?
    \item Czy agentowa pętla \textit{AutoResearch} przyspiesza dochodzenie do konfiguracji
    konkurencyjnych względem klasycznego podejścia eksperymentalnego?
\end{enumerate}

Dla powyższych pytań przyjęto hipotezy:
\begin{itemize}
    \item \textbf{H1}: konfiguracja YAML zmniejsza koszt odtwarzania eksperymentów
    i ogranicza liczbę błędów proceduralnych.
    \item \textbf{H2}: układ \textit{preprocessing + model + ensemble} istotnie wpływa
    na metryki per-era.
    \item \textbf{H3}: HPO podnosi jakość predykcji względem punktu odniesienia.
    \item \textbf{H4}: AutoResearch redukuje nakład manualnej iteracji przy zachowaniu
    porównywalnej lub lepszej jakości predykcji.
\end{itemize}

\section{Środowisko eksperymentalne}
\subsection{Dane i zakres}
Eksperymenty prowadzone są na danych turnieju Numerai Classic w wersji zgodnej z
implementacją projektu. Dane dzielone są zgodnie z logiką er i porządkiem czasowym,
z uwzględnieniem wymagań zapobiegania \textit{look-ahead bias}.

\subsection{Jednostka eksperymentu}
Za jednostkę eksperymentu przyjmuje się pojedyncze uruchomienie pipeline'u opisane
jednoznaczną konfiguracją (YAML lub konfiguracja wygenerowana przez HPO/AutoResearch),
które kończy się raportem metryk i zapisem artefaktów.

\subsection{Reprodukowalność}
Dla każdej próby zapisywane są:
\begin{itemize}
    \item konfiguracja wejściowa,
    \item ziarno losowe,
    \item metryki końcowe i pośrednie,
    \item identyfikatory artefaktów wyjściowych.
\end{itemize}
Zapewnia to odtwarzalność wyników i porównywalność między seriami eksperymentów.

\section{Metryki oceny}
\subsection{Metryki główne}
Ocena jakości opiera się na metrykach zgodnych z charakterem zadania Numerai:
\begin{itemize}
    \item średnia korelacja per-era,
    \item współczynnik Sharpe'a sygnału per-era.
\end{itemize}

\subsection{Metryki pomocnicze}
Dodatkowo analizowane są metryki wspierające interpretację:
\begin{itemize}
    \item odchylenie i rozrzut wyników między erami,
    \item stabilność wyników między uruchomieniami,
    \item koszt obliczeniowy (czas, liczba prób, zasoby).
\end{itemize}

\section{Projekt eksperymentów}
Plan badań podzielono na cztery etapy o rosnącej złożoności.

\subsection{Etap I: konfiguracja bazowa}
Celem etapu jest ustanowienie punktu odniesienia:
\begin{itemize}
    \item pojedynczy model bazowy,
    \item minimalny preprocessing,
    \item stałe parametry treningu i ewaluacji.
\end{itemize}
Wyniki etapu stanowią \textit{baseline} dla kolejnych porównań.

\subsection{Etap II: ablację komponentów}
W tym etapie badany jest wpływ poszczególnych elementów pipeline'u:
\begin{itemize}
    \item warianty preprocessingu,
    \item rodziny modeli bazowych,
    \item warianty ensemble (\texttt{single}, \texttt{weighted}, \texttt{stacking}).
\end{itemize}
Każda zmiana jest testowana przy możliwie stałych pozostałych warunkach.

\subsection{Etap III: optymalizacja hiperparametrów (HPO)}
Etap obejmuje uruchomienia Optuna dla wybranych rodzin modeli i konfiguracji:
\begin{itemize}
    \item definicja przestrzeni przeszukiwania,
    \item ustalenie budżetu prób,
    \item wybór najlepszych konfiguracji na podstawie metryki celu.
\end{itemize}
Rezultatem jest zestaw konfiguracji kandydackich do walidacji końcowej.

\subsection{Etap IV: pętla AutoResearch}
W ostatnim etapie uruchamiana jest agentowa pętla badawcza:
\begin{itemize}
    \item mutacje konfiguracji bazujące na historii wyników,
    \item iteracyjne testowanie propozycji,
    \item porównanie skuteczności względem HPO i baseline.
\end{itemize}
Ocena obejmuje zarówno jakość predykcji, jak i koszt dojścia do najlepszego wyniku.

\section{Protokół wykonania}
\subsection{Kontrola zmiennych}
Aby ograniczyć wpływ czynników zakłócających, stosowane są:
\begin{itemize}
    \item stałe podziały danych i procedury walidacji,
    \item kontrola ziaren losowych,
    \item jednolita ścieżka uruchomieniowa (CLI),
    \item ten sam sposób raportowania metryk i artefaktów.
\end{itemize}

\subsection{Kryteria zatrzymania}
Dla serii eksperymentalnych definiuje się:
\begin{itemize}
    \item maksymalny budżet prób/czasu,
    \item próg poprawy metryki (\textit{early stop} dla serii),
    \item minimalną liczbę uruchomień dla wnioskowania porównawczego.
\end{itemize}

\section{Metody analizy wyników}
\subsection{Analiza porównawcza}
Wyniki są analizowane na dwóch poziomach:
\begin{itemize}
    \item \textbf{poziom konfiguracji}: porównanie metryk między wariantami pipeline'u,
    \item \textbf{poziom procesu}: porównanie efektywności metod dochodzenia do
    dobrych konfiguracji (manualnie, HPO, AutoResearch).
\end{itemize}

\subsection{Analiza stabilności}
Dla najlepszych konfiguracji oceniana jest:
\begin{itemize}
    \item zmienność wyników między erami,
    \item odporność na losowość inicjalizacji,
    \item spójność jakości przy ponownym uruchomieniu.
\end{itemize}

\subsection{Interpretacja praktyczna}
Wnioski mają charakter nie tylko statystyczny, ale i inżynierski:
\begin{itemize}
    \item jaki wariant daje najlepszy kompromis jakość/koszt,
    \item które komponenty są krytyczne dla jakości,
    \item które etapy można efektywnie automatyzować.
\end{itemize}

\section{Zagrożenia dla trafności i ograniczenia metodyki}
W analizie uwzględnia się następujące ryzyka:
\begin{itemize}
    \item ograniczony budżet obliczeniowy zawężający przestrzeń testów,
    \item możliwe przeuczenie do konkretnego zakresu danych,
    \item zależność wyników od jakości zdefiniowanej przestrzeni HPO/AutoResearch.
\end{itemize}

Aby ograniczyć te ryzyka, plan zakłada:
\begin{itemize}
    \item eksperymenty ablation i walidację wielowariantową,
    \item jawne raportowanie konfiguracji i budżetów,
    \item oddzielenie wyników eksploracyjnych od końcowych wniosków porównawczych.
\end{itemize}

\section{Protokół reprodukowalności}
\label{sec:repro}

\subsection{Wersje środowiska}
\begin{itemize}
    \item Python: 3.12+,
    \item menedżer zależności: \texttt{uv},
    \item repozytorium: \todo{commit hash},
    \item dataset: Numerai Classic \datasetver{},
    \item pliki danych: \texttt{train.parquet}, \texttt{validation.parquet}, \texttt{features.json}.
\end{itemize}

\subsection{Komendy uruchomieniowe}
\begin{lstlisting}[basicstyle=\footnotesize\ttfamily,breaklines=true,breakatwhitespace=true,columns=fullflexible]
uv run python scripts/download_dataset.py --dataset-version v5.2 --output-dir data
uv run python scripts/run_experiment.py --config experiments/example_v1.yaml --artifact-dir artifacts/experiments
uv run python scripts/hpo_pipeline.py --data-dir data/v5.2 --train-subsample 0.125 --num-trials 30 --output-dir artifacts/hpo_x8 --local
uv run python scripts/autoresearch.py --data-dir data/v5.2 --train-subsample 0.125 --trials 50 --hours 2 --output-dir artifacts/autoresearch --agent-model claude-sonnet-4-6
uv run python scripts/export_numerai_pickle.py --data-dir data/v5.2 --best-config-path artifacts/hpo_x8/best_config.json --output-dir artifacts/competition_pickle
\end{lstlisting}

\subsection{Zasady raportowania}
Dla każdej próby zapisywane są: konfiguracja, seed, metryki per-era, czas wykonania, identyfikator artefaktu.

\section{Podsumowanie}
Przedstawiona metodyka łączy rygor badawczy z praktyką inżynierską: definiuje jasne
hipotezy, kontroluje warunki porównań i umożliwia ocenę zarówno jakości predykcji,
jak i efektywności procesu eksperymentalnego. W kolejnych rozdziałach metodyka ta
stanowi podstawę prezentacji wyników i ich krytycznej interpretacji.

\chapter{Wyniki eksperymentów i analiza}

\section{Cel rozdziału}
Celem rozdziału jest przedstawienie i interpretacja wyników uzyskanych z użyciem
frameworku \textit{AlphaPulse}. Analiza obejmuje porównanie konfiguracji bazowych,
wyniki ablacji komponentów, efekty optymalizacji hiperparametrów oraz ocenę skuteczności
pętli \textit{AutoResearch}. Rozdział odpowiada na pytania badawcze dotyczące jakości,
stabilności i efektywności procesu eksperymentalnego.

\section{Konfiguracja eksperymentów}
\subsection{Scenariusze porównawcze}
Eksperymenty zostały podzielone na cztery grupy:
\begin{itemize}
    \item \textbf{Baseline}: konfiguracje referencyjne bez zaawansowanej automatyzacji,
    \item \textbf{Ablacje}: kontrolowane zmiany preprocessingu, modeli i ensemble,
    \item \textbf{HPO}: przeszukiwanie hiperparametrów i konfiguracji pipeline'u,
    \item \textbf{AutoResearch}: iteracyjne mutacje konfiguracji proponowane agentowo.
\end{itemize}

\subsection{Zbierane artefakty}
Dla każdej próby rejestrowano:
\begin{itemize}
    \item konfigurację wejściową,
    \item metryki per-era i metryki zagregowane,
    \item czas wykonania i informacje o przebiegu uruchomienia,
    \item artefakty końcowe (m.in. najlepsze konfiguracje i predykcje).
\end{itemize}

\section{Wyniki konfiguracji bazowych}
\subsection{Punkt odniesienia}
Konfiguracje bazowe wyznaczyły poziom odniesienia dla dalszych porównań.
Pozwoliło to:
\begin{itemize}
    \item oszacować jakość sygnału możliwą bez intensywnej optymalizacji,
    \item zidentyfikować najbardziej wrażliwe elementy pipeline'u,
    \item określić realistyczny margines poprawy dla HPO i AutoResearch.
\end{itemize}

\subsection{Wnioski z baseline}
Wyniki bazowe potwierdziły, że nawet poprawna technicznie konfiguracja może być
niestabilna między erami, jeżeli nie uwzględnia właściwego doboru preprocessingu
i strategii łączenia modeli.

\section{Ablacja komponentów pipeline'u}
\subsection{Wpływ preprocessingu}
Analiza ablacyjna wykazała, że preprocessing istotnie wpływa na jakość i stabilność.
W szczególności:
\begin{itemize}
    \item metody skalowania poprawiały zachowanie części modeli liniowych i drzewiastych,
    \item selekcja/transformacja cech redukowała degradację na słabszych erach,
    \item nadmiernie agresywna redukcja wymiarowości mogła obniżać jakość sygnału.
\end{itemize}

\subsection{Wpływ doboru modeli}
Porównania rodzin modeli pokazały, że:
\begin{itemize}
    \item nie istnieje pojedynczy model dominujący we wszystkich ustawieniach,
    \item modele o wysokiej średniej jakości mogą różnić się stabilnością per-era,
    \item wartość modelu należy oceniać łącznie z preprocessingiem i ensemble.
\end{itemize}

\subsection{Wpływ strategii ensemble}
Strategie ensemble najczęściej poprawiały jakość względem pojedynczego modelu,
ale skala poprawy zależała od komplementarności modeli bazowych.
Najlepsze wyniki uzyskiwano, gdy:
\begin{itemize}
    \item modele wnosiły zróżnicowany sygnał,
    \item agregacja była dostrojona do metryki celu,
    \item unikało się nadmiernego dopasowania meta-warstwy.
\end{itemize}

\section{Wyniki HPO}
\subsection{Skuteczność optymalizacji}
HPO systematycznie poprawiało wyniki względem baseline dla większości badanych
konfiguracji. Najczęstsze źródła poprawy to:
\begin{itemize}
    \item lepsze dopasowanie hiperparametrów modelu bazowego,
    \item korzystniejsze ustawienia preprocessingu,
    \item bardziej trafny dobór parametrów ensemble.
\end{itemize}

\subsection{Koszt optymalizacji}
Poprawa jakości wymagała istotnego budżetu obliczeniowego.
Zaobserwowano klasyczny kompromis:
\begin{itemize}
    \item większa liczba prób zwiększa szansę na znalezienie lepszego punktu,
    \item krańcowy zysk jakości maleje wraz z kolejnymi próbami.
\end{itemize}

\subsection{Charakter najlepszych konfiguracji}
Najlepsze konfiguracje zwykle nie były ekstremalne, lecz \textit{zbalansowane}:
umiarkowana złożoność pipeline'u, stabilny preprocessing i rozsądnie dobrana
agregacja predykcji.

\section{Wyniki AutoResearch}
\subsection{Efektywność pętli agentowej}
AutoResearch pozwolił zautomatyzować iteracyjne generowanie hipotez konfiguracyjnych.
W praktyce:
\begin{itemize}
    \item zmniejszał liczbę manualnych decyzji eksperymentalnych,
    \item przyspieszał przejście od konfiguracji bazowych do kandydatów wysokiej jakości,
    \item wspierał eksplorację wariantów pomijanych w ręcznym strojeniu.
\end{itemize}

\subsection{Porównanie z HPO}
AutoResearch i HPO mają częściowo komplementarny charakter:
\begin{itemize}
    \item HPO jest silniejsze w lokalnym dostrajaniu parametrów,
    \item AutoResearch sprawdza się w szerszej eksploracji zmian strukturalnych.
\end{itemize}
Najbardziej efektywne okazało się podejście hybrydowe:
eksploracja agentowa, a następnie doprecyzowanie przez HPO.

\section{Analiza stabilności i ryzyka overfittingu}
\subsection{Stabilność per-era}
Ocena per-era wykazała, że wysoka średnia metryka nie gwarantuje stabilności.
Dlatego konfiguracje finalne wybierano z uwzględnieniem:
\begin{itemize}
    \item rozrzutu wyników między erami,
    \item odporności na słabsze okresy,
    \item powtarzalności przy ponownym uruchomieniu.
\end{itemize}

\subsection{Kontrola przeuczenia}
Aby ograniczyć overfitting:
\begin{itemize}
    \item stosowano walidację zgodną z charakterem czasowym danych,
    \item analizowano różnice między jakością lokalną i zagregowaną,
    \item preferowano konfiguracje stabilne nad niestabilnie „rekordowe”.
\end{itemize}

\section{Syntetyczne odpowiedzi na pytania badawcze}
\begin{itemize}
    \item \textbf{P1 (reprodukowalność):} podejście deklaratywne i artefaktowe zwiększyło
    porównywalność eksperymentów i ułatwiło odtwarzanie wyników.
    \item \textbf{P2 (wpływ komponentów):} preprocessing, model bazowy i ensemble mają
    istotny, wzajemnie zależny wpływ na metryki per-era.
    \item \textbf{P3 (rola HPO):} HPO poprawia wyniki względem baseline, ale wymaga
    kontrolowanego budżetu obliczeniowego.
    \item \textbf{P4 (rola AutoResearch):} pętla agentowa redukuje koszt manualnej iteracji
    i skutecznie wspiera eksplorację przestrzeni konfiguracji.
\end{itemize}

\section{Ograniczenia interpretacji wyników}
Interpretacja wyników wymaga uwzględnienia:
\begin{itemize}
    \item ograniczonego budżetu eksperymentalnego,
    \item zależności od konkretnej wersji danych i zakresu testów,
    \item wpływu doboru przestrzeni przeszukiwania na wyniki HPO i AutoResearch.
\end{itemize}

Ograniczenia te nie podważają wniosków kierunkowych, lecz wyznaczają granice ich
uogólniania.

\section{Agregacja wyników liczbowych}
\label{sec:wyniki-liczbowe}

\subsection{Porównanie scenariuszy}
\begin{table}[h]
\centering
\caption{Baseline vs HPO vs AutoResearch (validation)}
\begin{tabular}{lcccc}
\toprule
Scenariusz & Mean corr/era & Sharpe/era & Czas [h] & Liczba prób \\
\midrule
Baseline & \todo{0.0XX} & \todo{0.XX} & \todo{0.X} & 1 \\
HPO & \todo{0.0XX} & \todo{0.XX} & \todo{X.X} & \todo{30} \\
AutoResearch & \todo{0.0XX} & \todo{0.XX} & \todo{X.X} & \todo{50} \\
\bottomrule
\end{tabular}
\end{table}

\subsection{Top konfiguracje}
\begin{table}[h]
\centering
\caption{Najlepsze konfiguracje końcowe}
\begin{tabular}{clccc}
\toprule
Rank & Konfiguracja (model + ensemble) & Corr/era & Sharpe/era & Stabilność \\
\midrule
1 & \todo{XGBoost + weighted} & \todo{...} & \todo{...} & \todo{...} \\
2 & \todo{LightGBM + stacking(ridge)} & \todo{...} & \todo{...} & \todo{...} \\
3 & \todo{CatBoost + single} & \todo{...} & \todo{...} & \todo{...} \\
\bottomrule
\end{tabular}
\end{table}

\subsection{Wyniki negatywne}
\begin{itemize}
    \item \todo{np. agresywne PCA obniżało corr o ...}
    \item \todo{np. stacking bez regularyzacji meta-modelu destabilizował wyniki}
    \item \todo{np. TabPFN/TabICL: ograniczenia kosztu/czasu przy pełnym subsample}
\end{itemize}

\section{Podsumowanie}
Wyniki potwierdzają, że \textit{AlphaPulse} umożliwia systematyczne i reprodukowalne
eksperymentowanie z potokami predykcyjnymi dla Numerai. Największą wartość praktyczną
stanowi połączenie modularnego pipeline'u, HPO i pętli AutoResearch, które razem
zwiększają jakość predykcji oraz efektywność procesu badawczego.

\chapter{Dyskusja, wnioski i kierunki dalszych prac}

\section{Cel rozdziału}
Celem rozdziału jest syntetyczne podsumowanie rezultatów pracy, krytyczna dyskusja
uzyskanych wyników oraz wskazanie kierunków dalszego rozwoju frameworku
\textit{AlphaPulse}. Rozdział domyka odpowiedzi na pytania badawcze i formułuje
praktyczne rekomendacje dla kolejnych iteracji systemu.

\section{Dyskusja wyników}
\subsection{Znaczenie uzyskanych efektów}
Przeprowadzone badania pokazują, że największa wartość projektu nie wynika z pojedynczego
modelu, lecz z podejścia systemowego: połączenia konfigurowalnej architektury,
modularnego pipeline'u, automatyzacji strojenia i agentowej eksploracji konfiguracji.

W praktyce oznacza to, że:
\begin{itemize}
    \item proces badawczy staje się bardziej powtarzalny,
    \item porównania wariantów są metodycznie uporządkowane,
    \item poprawa wyników nie jest efektem przypadkowych zmian, lecz kontrolowanych iteracji.
\end{itemize}

\subsection{Jakość predykcji a stabilność}
Jednym z kluczowych wniosków jest konieczność równoważenia jakości średniej i stabilności
per-era. Konfiguracje osiągające lokalnie najwyższe metryki nie zawsze zapewniały
najlepszą odporność na zmienność danych.

Z perspektywy zastosowania konkursowego bardziej użyteczne są konfiguracje:
\begin{itemize}
    \item o nieco niższej jakości szczytowej,
    \item ale większej spójności wyników między erami,
    \item i lepszej powtarzalności przy ponownych uruchomieniach.
\end{itemize}

\subsection{Rola automatyzacji}
Wyniki potwierdzają, że automatyzacja jest niezbędna przy dużej przestrzeni konfiguracji.
HPO i AutoResearch pełnią różne role:
\begin{itemize}
    \item HPO skutecznie dostraja parametry lokalnie,
    \item AutoResearch ułatwia eksplorację strukturalnych zmian pipeline'u.
\end{itemize}
Najbardziej efektywne okazuje się podejście łączone, w którym agent proponuje kierunki
zmian, a HPO je doprecyzowuje.

\section{Wnioski końcowe}
\subsection{Wnioski względem celu głównego}
Cel pracy, polegający na zaprojektowaniu i implementacji konfigurowalnego środowiska
AutoML dla finansowych potoków predykcyjnych Numerai, został osiągnięty.
Opracowany framework:
\begin{itemize}
    \item integruje pełny przepływ od danych do artefaktu submisji,
    \item umożliwia deklaratywne definiowanie eksperymentów,
    \item wspiera systematyczne i reprodukowalne badania konfiguracji modeli.
\end{itemize}

\subsection{Wnioski względem pytań badawczych}
\begin{itemize}
    \item \textbf{Abstrakcje komponentów}: modularny podział na preprocessing, modele
    i ensemble jest wystarczająco elastyczny dla szerokiego zakresu konfiguracji.
    \item \textbf{Walidacja}: uwzględnienie struktury czasowej i er jest krytyczne dla
    wiarygodności wyników.
    \item \textbf{Automatyzacja}: HPO oraz AutoResearch istotnie wspierają proces
    poszukiwania konfiguracji konkurencyjnych jakościowo.
    \item \textbf{Integracja workflow}: system zapewnia spójny przepływ eksperymentalny
    od konfiguracji do eksportu predykcji.
\end{itemize}

\subsection{Wkład praktyczny i naukowy}
Wkład praktyczny pracy obejmuje gotowy do użycia framework eksperymentalny.
Wkład naukowo-metodyczny polega na wykazaniu, że w zadaniu typu Numerai przewagę
daje dobrze zaprojektowany proces badawczy, a nie tylko jednorazowy dobór modelu.

\section{Ograniczenia pracy}
Pomimo uzyskanych rezultatów, praca ma ograniczenia:
\begin{itemize}
    \item eksperymenty realizowano w określonym budżecie obliczeniowym,
    \item zakres badań dotyczy konkretnej wersji danych i środowiska konkursowego,
    \item skuteczność automatyzacji zależy od jakości zdefiniowanej przestrzeni przeszukiwania.
\end{itemize}

Ograniczenia te wyznaczają naturalny kierunek dalszych prac i nie podważają głównych
wniosków dotyczących architektury i metodyki.

\section{Kierunki dalszego rozwoju}
\subsection{Rozwój metod ewaluacji}
W kolejnych iteracjach warto rozszerzyć system metryk o:
\begin{itemize}
    \item dodatkowe miary zgodności z celami payout Numerai,
    \item bardziej szczegółowe miary stabilności sygnału w czasie,
    \item analizy niepewności i odporności na przesunięcia rozkładów danych.
\end{itemize}

\subsection{Rozwój automatyzacji}
Obiecującym kierunkiem jest dalsza integracja warstwy agentowej:
\begin{itemize}
    \item adaptacyjne zarządzanie budżetem eksperymentalnym,
    \item inteligentne zawężanie przestrzeni HPO na podstawie historii prób,
    \item automatyczne raporty porównawcze konfiguracji.
\end{itemize}

\subsection{Rozwój architektury systemu}
W kontekście inżynierskim zasadne są:
\begin{itemize}
    \item dalsza standaryzacja interfejsów komponentów,
    \item rozszerzenie katalogu modeli i preprocessingu,
    \item lepsze mechanizmy walidacji konfiguracji i diagnostyki błędów.
\end{itemize}

\subsection{Rozszerzenie zastosowań}
Choć system był projektowany pod Numerai Classic, architektura ma potencjał adaptacji
do innych zadań tablicowych i finansowych, pod warunkiem dostosowania warstwy danych,
metryk i procedur walidacyjnych.

\section{Podsumowanie rozdziału}
Praca potwierdziła, że skuteczne modelowanie w środowisku finansowym wymaga jednocześnie:
poprawnej metodologii walidacyjnej, modularnej architektury oprogramowania oraz
zautomatyzowanego procesu eksperymentalnego. \textit{AlphaPulse} spełnia te warunki,
stanowiąc solidną podstawę do dalszych badań i rozwoju praktycznych rozwiązań AutoML
dla zadań predykcyjnych o wysokiej złożoności.

\section{Odpowiedzi liczbowe na pytania badawcze}
\begin{itemize}
    \item \textbf{P1 (reprodukowalność):} \todo{np. 100\% prób odtworzonych z tej samej konfiguracji}.
    \item \textbf{P2 (wpływ komponentów):} największy wkład: \todo{ensemble/preprocessing/model X}.
    \item \textbf{P3 (HPO):} poprawa względem baseline: \todo{+X.XX corr, +Y Sharpe}.
    \item \textbf{P4 (AutoResearch):} redukcja czasu iteracji: \todo{-Z\%}, koszt prób: \todo{N}.
\end{itemize}

\chapter{Bibliografia}
Bibliografia [1] Palamabron/AlphaPulse (feat/core-framework). \url{https://github.com/Palamabron/AlphaPulse/tree/feat/core-framework}
[2] Numerai tournament: Overview (docs). \url{https://docs.numer.ai/numerai-tournament/readme.md}
[3] López de Prado, Marcos. Advances in Financial Machine Learning. Wiley, 2018. \url{https://www.wiley.com/en-us/Advances+in+Financial+Machine+Learning-p-9781119482086}
[4] eslazarev. Purged-cross-validation (library: purging + embargo). \url{https://github.com/eslazarev/purged-cross-validation}
[5] Bailey, David H., and López de Prado, Marcos. The Deflated Sharpe Ratio: Correcting for Selection Bias, Backtest Overfitting and Non-Normality. \url{https://davidhbailey.com/dhbpapers/deflated-sharpe.pdf}
[6] Akiba, T., Sano, S., Yanase, T., Ohta, T., and Koyama, M. Optuna: A Next-generation Hyperparameter Optimization Framework. \url{https://arxiv.org/pdf/1907.10902}
[7] He, Xin, Zhao, Kaiyong, and Chu, Xiaowen. AutoML: A Survey of the State-of-the-Art. Knowledge-Based Systems, 2021. \url{https://arxiv.org/abs/1908.00709}
[8] Ferrante, A., et al. Benchmark and Survey of Automated Machine Learning Frameworks. Journal of Artificial Intelligence Research, 2021. \url{https://pdfs.semanticscholar.org/b98d/cbe6c7e14827f280a164c08ae5a63005bb7c.pdf}
[9] Chen, Tianqi, and Guestrin, Carlos. XGBoost: A Scalable Tree Boosting System. \url{https://arxiv.org/pdf/1603.02754}
[10] Ke, Guolin, Meng, Qi, Finley, Thomas, et al. LightGBM: A Highly Efficient Gradient Boosting Decision Tree. \url{https://proceedings.neurips.cc/paper/2017/file/6449f44a102fde848669bdd9eb6b76fa-Paper.pdf}
[11] Prokhorenkova, Liudmila, Gusev, Gleb, Vorobev, Alexander, et al. CatBoost: Unbiased Boosting with Categorical Features. \url{https://proceedings.neurips.cc/paper_files/paper/2018/file/14491b756b3a51daac41c24863285549-Paper.pdf}
[12] Wolpert, David H. Stacked Generalization. Neural Networks, 1992. \url{https://ui.adsabs.harvard.edu/abs/1992NN......5..241W/abstract}
[13] Wolpert, David H. (stacking) + meta-level issues: Issues in Stacked Generalization. \url{https://jair.org/index.php/jair/article/download/10228/24333}
[14] Breiman, Leo. Bagging Predictors. Machine Learning, 1996. \url{https://sci2s.ugr.es/keel/pdf/algorithm/articulo/1996-ML-Breiman-Bagging%20Predictors.pdf}
[15] Hollmann, Nathan, et al. TabPFN: A Transformer That Solves Small Tabular Classification Problems in a Second. \url{https://arxiv.org/pdf/2207.01848}
[16] Hollmann, Nathan, et al. Accurate predictions on small data with a tabular foundation model (TabPFN). Nature, 2024. \url{https://preview-www.nature.com/articles/s41586-024-08328-6.pdf}
\end{document}
